\scanpage{69}
\mspage{68}
% A long oblique cancellation crosses the body of this draft, reaching the final paragraph.

\deleted{Ceci peut se généraliser ainsi. Soit}
\[
\deleted{P\subset\unsure{L|\mathcal{L}}}
\]
\margin{$P$ comme ``points''}
\deleted{une partie de $L$ (ou simplement de $\unsure{\mathcal{F}|\mathscr{F}}$ --- mais pratiquement ce sera tjrs $\subset L$ je présume) telle que}

\deleted{a) si $x,y\in P$, $x\ne y\Rightarrow x\parallel y$ (donc sur $P$, $x\parallel y$ équivaut à $x\ne y$)}

\deleted{b) si $F,G\in\mathcal{F}$ (resp. $M\unsure{\}}$, si on part de $M$) \unsure{et}}

\deleted{\unsure{alors} \unsure{et si} $\forall x\in\operatorname{omb}_{\unsure{\Lambda|P}}(F),\ y\in\operatorname{omb}_{\unsure{\Lambda|P}}(G)$, on a}

\[
\deleted{\unsure{F\parallel G}}
\]
\deleted{On a $F\parallel G$ si (\unsure{et s.s.|et ss.}) $\nexists x\in P$, $\unsure{\operatorname{supp}x\subset\operatorname{supp}F}$}
\inserted{\deleted{avec $\operatorname{supp}x\subset\operatorname{supp}A\cap\operatorname{supp}B$}}

\deleted{\unsure{si} $x\parallel y$, alors $F\parallel G$}

\deleted{(C'est la condition $\unsure{\mathrm{AtP2}|\mathrm{At}\Lambda2}$, \unsure{i.e.|ou} $\mathrm{At}\Lambda2$ dans le cas $\Lambda=P$ --- condition qui implique $\mathrm{AtL2}$.)}

% margin note written slanted; its arrow points to the paragraph beginning Alors.
\margin{Un peu bref ! $\downarrow$}
\deleted{Alors l'ens ordonné des supports s'identifie à}
\[
\deleted{\Sigma_P\simeq\unsure{\mathscr{P}|\mathcal{P}}(P),\quad\text{donc}}
\]
\diagram{\begin{tikzcd}[ampersand replacement=\&, column sep=large, execute at end picture={\draw ([xshift=-1mm]current bounding box.north west)--([xshift=1mm]current bounding box.south east);}]
\deleted{\unsure{s}}\quad\deleted{A} \arrow[r, mapsto] \& \deleted{\{x\in P\mid\operatorname{supp}x\subset A\}} \\
\deleted{\Sigma_{\deleted{\unsure{\mathcal{F}}}\unsure{F|\mathcal{F}}\text{ (ou }M\text{, ou }L\text{)}}} \arrow[r, phantom, "\deleted{\simeq}"] \& \deleted{\mathscr{P}(P)}
\end{tikzcd}}
\deleted{iso d'ens. ordonnés, compatible à l'application $C$ (``complémentaire''). Donc dans ce cas encore, les perplexités \unsure{s'évanouissent}. Dans ce cas, il y a lieu de définir le support comme}
\[
\deleted{\operatorname{supp}A=\operatorname{omb}_\Lambda A=\bigcup_{F\in A}\operatorname{omb}_\Lambda(F).}
\]

\deleted{Notons que l'isom. d'ens. ordonnés avec \unsure{opération} ``passage au complémentaire''}
\[
\deleted{\Sigma\simeq\mathscr{P}(P)}
\]
